\section*{Hoofdstuk \ref{ch:g9:alternating-current} --- Wisselstroom}

\begin{solution}{exo:g9:alternating-current:1}
\omterm{def:g9:alternating-current:dcac}{Gelijkstroom} loopt altijd één kant op onder een gestage
\omterm{def:g8:voltage-voltmeter:voltage}{spanning} --- het regime van de \omterm{def:g3:first-electric-circuit:battery}{batterij}.
\omterm{def:g9:alternating-current:dcac}{Wisselstroom} keert periodiek onder een zwaaiende
\omterm{def:g8:voltage-voltmeter:voltage}{spanning} --- het regime van het stopcontact.
\end{solution}

\begin{solution}{exo:g9:alternating-current:2}
$T$: de duur van één volledige kringloop; $f$: kringlopen per
\omterm{def:g2:measuring-time:second}{seconde}; $f = 1/T$ en $T = 1/f$. Voor \qty{50}{Hz}: $T =
\qty{0.02}{s}$. Voor $T = \qty{0.01}{s}$: $f = \qty{100}{Hz}$.
\end{solution}

\begin{solution}{exo:g9:alternating-current:3}
Honderd keer: twee omkeringen in elk van de vijftig kringlopen.
\end{solution}

\begin{solution}{exo:g9:alternating-current:4}
$T = 4 \times 5 = \qty{20}{ms} = \qty{0.02}{s}$; $f = 1/0.02 =
\qty{50}{Hz}$.
\end{solution}

\begin{solution}{exo:g9:alternating-current:5}
\qty{230}{V} is de \omterm{prop:g9:alternating-current:effective}{effectieve spanning} --- de gestage
gelijkstroomwaarde die even goed zou verwarmen, en wat
wisselstroomvoltmeters aflezen; \qty{325}{V} is de piek, de
top die de zwaai werkelijk aanraakt, beide kanten op, vijftig keer per
\omterm{def:g2:measuring-time:second}{seconde}.
\end{solution}

\begin{solution}{exo:g9:alternating-current:6}
Verwarmen geeft niets om richting --- de spiraal wordt warm van elke
mars, dus drinken waterkokers de zwaai rechtstreeks. Laden is
scheikunde achterstevoren, en scheikunde eist één vaste richting: het
blok zet de zwaai eerst om in standvastigheid.
\end{solution}

\begin{solution}{exo:g9:alternating-current:7}
De \omterm{def:g3:first-electric-circuit:battery}{batterij}: een vlakke lijn op \qty{9}{V}. De wisselspanning:
een soepele golf die tussen $+9$ en \qty{-9}{V} zwaait, één kringloop
per \qty{20}{ms}. De oogopslagtoets: vlak tegenover golvend.
\end{solution}

\begin{solution}{exo:g9:alternating-current:8}
Draaiende machines maken van nature \omterm{def:g9:alternating-current:dcac}{wisselstroom}; en alleen
\omterm{def:g9:alternating-current:dcac}{wisselstroom} voedt de transformator, zonder welke elektriciteit het
land niet op hoge \omterm{def:g8:voltage-voltmeter:voltage}{spanning} zou kunnen doorkruisen en huizen niet
op een tamme zou kunnen binnenkomen.
\end{solution}

\begin{solution}{exo:g9:alternating-current:9}
Bij langzaam trappen laten de weinige kringlopen per \omterm{def:g2:measuring-time:second}{seconde} van de
dynamo het \omterm{def:g3:first-electric-circuit:bulb}{lampje} zichtbaar pulseren --- elke top een flits, elke
nuldoorgang een dip; sneller trappen vermenigvuldigt de pulsen tot het
oog ze samensmelt. De $24$ sporen van de bioscoop leggen de ruwe
samensmeltdrempel van het oog vast: onder ongeveer \qty{25}{Hz} aan
lichtpulsen stoort het geflikker; erboven leest de gloed als gestaag.
\end{solution}

\begin{solution}{exo:g9:alternating-current:10}
Toppen op ongeveer \qty{325}{V} --- en honderd per \omterm{def:g2:measuring-time:second}{seconde}, vijftig
positieve en vijftig negatieve.
\end{solution}

\begin{solution}{exo:g9:alternating-current:11}
$T = 1/60 \approx \qty{0.017}{s}$; piek $\approx 120 \times 1.4 \approx
\qty{170}{V}$. Kachels en lampen passen zich onverschillig aan; moderne
laders lezen ``100--240 V, 50/60 Hz'' op hun etiket en zetten vrolijk
om --- de bezwaarmakers zijn oude apparaten voor één enkele
\omterm{def:g8:voltage-voltmeter:voltage}{spanning}, en elke motorklok die de kringlopen van het net
telt: die zou op zestig te snel lopen.
\end{solution}

\begin{solution}{exo:g9:alternating-current:12}
Oplichten bij beide toppen geeft $2 \times 50 = 100$ pulsen per
\omterm{def:g2:measuring-time:second}{seconde}. Een camera van $25$ sporen bemonstert elke \qty{0.04}{s} ---
vier pulsen per spoor, maar niet op elk punt van de rollende belichting
gelijk uitgelijnd: het kleine verschil tussen het ritme van de camera
en dat van de lamp schildert banden die over opeenvolgende sporen
kruipen --- een zweving tussen twee klokken. Bij \qty{60}{Hz} ($120$
pulsen) verandert het verschil met de $25$ sporen, en veranderen de
balken van breedte en \omterm{def:g7:motion-average-speed:formula}{snelheid} --- het patroon verraadt het
plaatselijke net.
\end{solution}

\begin{solution}{pb:g9:alternating-current:1}
\textbf{1.} \omterm{def:g9:alternating-current:dcac}{Gelijkstroom} op $2.25 \times 2 = \qty{4.5}{V}$: de platte
zakbatterij.
\textbf{2.} \omterm{def:g9:alternating-current:dcac}{Wisselstroom}; $T = \qty{20}{ms}$, $f = \qty{50}{Hz}$; piek
$3 \times 2 = \qty{6}{V}$.
\textbf{3.} \omterm{def:g9:alternating-current:dcac}{Gelijkstroom} van \qty{4.5}{V} \emph{omgekeerd} aangesloten
--- de snoeren verwisseld: onschuldig op een \omterm{met:g9:alternating-current:scope}{oscilloscoop}, leerzaam in
het grootboek.
\textbf{4.} \omterm{def:g9:alternating-current:dcac}{Wisselstroom}; $T = \qty{10}{ms}$, $f = \qty{100}{Hz}$; piek
\qty{3}{V}. Haar honderd verdubbelde pulsen per \omterm{def:g2:measuring-time:second}{seconde} liggen
comfortabeler voorbij de samensmeltdrempel van het oog dan het
geflikker van vijftig kringlopen van B.
\textbf{5.} De \omterm{def:g8:voltage-voltmeter:voltmeter}{voltmeter} meldt de effectieve waarde --- het
gelijkstroomequivalent voor verwarming --- die voor een sinusoïde
ongeveer op de piek gedeeld door $1.4$ ligt: $6 \div 1.4 \approx
\qty{4.3}{V}$. Beide getallen beschrijven één zwaai.
\textbf{6.} Alleen voeding A (goed om --- en C nadat hij is
omgestoken). De wisselstroomvoedingen hebben een omzetter nodig --- het
gelijkrichtende blok van de lader --- tussen henzelf en welke
\omterm{def:g3:first-electric-circuit:battery}{batterij} dan ook.
\textbf{7.} De vijftig zwaaien per \omterm{def:g2:measuring-time:second}{seconde} van het lichtnet liggen
dicht bij het eigen elektrische ritme van het hart en kunnen het
ontregelen: \omterm{def:g8:intensity-ammeter:intensity}{ampère} voor \omterm{def:g8:intensity-ammeter:intensity}{ampère} is de zwaai
gevaarlijker dan de gestage mars.
\textbf{8.} De sinusoïde --- de volmaakte zwaai, van nature getekend
door een spoel die met constante \omterm{def:g7:motion-average-speed:formula}{snelheid} in een magnetisch
veld draait: de wisselstroomgenerator van het volgende hoofdstuk.
\textbf{9.} Eén kringloop per $4$ hokjes (\qty{20}{ms}); toppen $325
\div 100 = 3.25$ hokjes boven en onder.
\textbf{10.} $f = 1/0.02 = \qty{50}{Hz}$; effectief $\approx 325 \div
1.4 \approx \qty{230}{V}$ --- het etiket, teruggewonnen van het scherm.
\textbf{11.} Een vlakke lijn op \qty{230}{V} --- de standvastigheid van
een gelijk verwarmingsvermogen.
\textbf{12.} Bijvoorbeeld: ``Een vlakke lijn is de standvastigheid van
een \omterm{def:g3:first-electric-circuit:battery}{batterij} --- één richting, één waarde. Een golf is de zwaai
van het net --- richting en waarde in een onophoudelijke kringloop. En
de eerlijke wisselkoers tussen die twee is de
\omterm{prop:g9:alternating-current:effective}{effectieve spanning}: gelijke \omterm{def:g5:heat-and-insulation:heat}{warmte}, gelijke
waarde.''
\end{solution}
